\documentclass[10pt,a4paper]{amsart}
\usepackage[margin=19mm]{geometry}
\usepackage{amsmath,amssymb,mathtools}
\usepackage[scr=rsfs,cal=euler]{mathalfa}
\usepackage{xfrac}
\usepackage[unicode,hidelinks,bookmarks=false]{hyperref}
\newcommand{\ie}{i.e.\ }
\newcommand{\cf}{cf.\ }
\newcommand{\resp}{respectively\ }
\newcommand{\Subsection}{\subsection}
\providecommand{\nobreakdash}{\nobreak\hbox{-}}
\setlength{\emergencystretch}{2em}
\multlinegap0pt
\usepackage{amssymb}
\usepackage{enumerate}
\newtheorem{lem}{Lemma}
\newcommand{\R}{\mathbb{R}}
\title{The Dirichlet problem for the prescribed scalar curvature in Anti-de Sitter space}
\author{Pierre Bayard}
\date{Source: arXiv:2606.07360v1 (CC BY 4.0)}
\begin{document}
\maketitle

\noindent\emph{Source and licence.} The Dirichlet problem for the prescribed scalar curvature in Anti-de Sitter space. Version arXiv:2606.07360v1 (CC BY 4.0), distributed by arXiv under the Creative Commons Attribution 4.0 International licence, http://creativecommons.org/licenses/by/4.0/, as recorded in the arXiv licence metadata for this exact version and shown on the abstract page https://arxiv.org/abs/2606.07360v1. Full licence notice: \texttt{licenses/arxiv-2606.07360-CC-BY-4.0.txt}.\par\smallskip

\noindent\textit{Editorial scope.} Counted unit: the article's lem nablamunu (source0.tex lines 345--351). The article's complete original proof of it is delivered with it (source0.tex lines 352--391, one \texttt{proof} environment of 40 lines). No mathematics is written, repaired or re-derived: the statement and the proof are the article's own lines, in the article's own notation, and no retained line is substituted or re-flowed.  Retained uncounted as the source-local context the proof needs: lines 305--307; lines 309--311; lines 330--332; lines 334--336.  Nothing was omitted from any retained span, so the package carries no omission record.  No retained line contains a citation, so the wrapper carries no bibliography and no bibliography entry.  No retained line contains a figure environment, an image-inclusion command or drawing code, so no image is delivered and the package declares no asset dependency.  Editorial formatting only, in addition: an AMS \texttt{amsart} preamble that restates the article's own declarations and macros used by the retained lines, namely the environments \texttt{lem} on separate counters, together with the standard packages those lines need.  The retained environments print in this document as Lemma 1 in order of appearance, whereas the article prints the same items with the numbering of its own full document.  The complete native source of the article as distributed in the arXiv e-print for this exact version is bundled unchanged as \texttt{sources/202253/source0.tex}.
\begin{equation}\label{notation est C1 expr nu}
\nu=\frac{1}{\mu}\frac{1}{\sqrt{1-|\nabla^Su|^2}}
\end{equation}

\begin{equation}\label{notation est C1 expr N}
N=\nu\left(\nabla^Su+e_{n+1}\right).
\end{equation}

\begin{equation}\label{C1 christoffel 1}
\langle\widetilde{\nabla}_Xe_{\alpha},e_{n+1}\rangle=-X^{n+1}\ \mu\ e_\alpha.\mu=-X.u\ \mu\ e_\alpha.\mu,\ \alpha=1,\ldots,n,
\end{equation}

\begin{equation}\label{C1 christoffel 2}
\langle\widetilde{\nabla}_Xe_{n+1},e_{n+1}\rangle=-\mu\ X.\mu,
\end{equation}

\begin{lem}\label{lem nablamunu}
The gradient of the function $\nu$ is given by
\begin{equation}\label{nablamunu}
\nabla(\mu\nu)=-\frac{1}{\nu}\langle\nabla\mu,\nabla u\rangle\nabla u+\mu S(\nabla u)
\end{equation}
where $S:TM\rightarrow TM$ is the shape operator of $M.$
\end{lem}

\begin{proof}
Since $\nu=\langle T,N\rangle,$ we have
\begin{equation}\label{eqn dnu pf nablamunu}
d\nu(X)=\langle\widetilde{\nabla}_XT,N\rangle+\langle T,\widetilde{\nabla}_XN\rangle
\end{equation}
for all $X\in TM$. We compute the first term on the right-hand side of (\ref{eqn dnu pf nablamunu}): since $T=-1/\mu^2\ e_{n+1}$ we easily get using (\ref{C1 christoffel 1}) and (\ref{C1 christoffel 2}) that
$$\widetilde{\nabla}_XT=-\frac{1}{\mu}X.u\ \nabla^H\mu+\frac{1}{\mu^3}X.\mu\ e_{n+1}$$
and using (\ref{notation est C1 expr N}) with $\nabla^Su=\mu^2\nabla^Hu$ we deduce that
\begin{eqnarray}
\langle\widetilde{\nabla}_XT,N\rangle&=&-\mu\nu X.u\ \langle\nabla^H\mu,\nabla^H u\rangle-\frac{\nu}{\mu}X.\mu\nonumber\\
&=&-\frac{1}{\mu\nu} X.u\ \langle\nabla\mu,\nabla u\rangle-\frac{\nu}{\mu}X.\mu.\label{dnu interm1}
\end{eqnarray}
We used in the last step the formula
\begin{equation}\label{pf lem grad nu eqn interm} 
\langle\nabla\mu,\nabla u\rangle=\mu^2\nu^2\langle\nabla^H\mu,\nabla^H u\rangle;
\end{equation}
on the left-hand side $\nabla\mu$ and $\nabla u$ belong to $TM$, and on the right-hand side $\nabla^H\mu$ and $\nabla^H u$ belong to $TB^n$ and are gradients computed with respect to the hyperbolic metric $g_H.$ We will prove that formula below, but we first compute the second term of the right-hand side of (\ref{eqn dnu pf nablamunu}) and conclude the proof of (\ref{nablamunu}): noticing that the orthogonal projection of $T$ on $TM$ is $T^{||}=\nabla u,$ we have
\begin{equation}\label{dnu interm2}
\langle T,\widetilde{\nabla}_XN\rangle=\langle T^{||},S(X)\rangle=\langle\nabla u,S(X)\rangle=\langle S(\nabla u),X\rangle,
\end{equation}
and (\ref{eqn dnu pf nablamunu}), together with (\ref{dnu interm1}) and (\ref{dnu interm2}), yields
$$\nabla\nu=-\frac{1}{\mu\nu}\langle\nabla\mu,\nabla u\rangle\nabla u-\frac{\nu}{\mu}\nabla\mu+S(\nabla u),$$
which gives (\ref{nablamunu}). We finally prove (\ref{pf lem grad nu eqn interm}). We first note that a vector $X\in TM$ is of the form $X=X'+\langle\nabla^Hu,X'\rangle\ e_{n+1}$ with $X'\in TB^n$. By composition with the natural projection $M\rightarrow B^n,$ a function $f:B^n\rightarrow\R$ may be considered as a function $f:M\rightarrow\R,$ and by the observation above its gradient may be written $\nabla f=(\nabla f)'+\langle\nabla ^H u,(\nabla f)'\rangle\ e_{n+1}$ where $(\nabla f)'$ belongs to $TB^n.$ We have, for all $X=X'+\langle\nabla^Hu,X'\rangle\ e_{n+1}\in TM,$
$$df(X')=df(X)=\langle\nabla f,X\rangle=\langle (\nabla f)',X'\rangle-\mu^2\langle\nabla ^H u,(\nabla f)'\rangle\langle\nabla^Hu,X'\rangle$$
which implies that 
$$\nabla^Hf=(\nabla f)'-\mu^2\langle\nabla^Hu,(\nabla f)'\rangle\nabla^H u.$$ 
Taking the scalar product with $\nabla^Hu$ and using that 
\begin{equation}\label{est C1 expr gradH mu}
\mu^2|\nabla^Hu|^2=1-1/{\mu^2\nu^2}
\end{equation} 
(as a consequence of (\ref{notation est C1 expr nu}), using that $\nabla^Su=\mu^2\nabla^Hu$ and where the norm of $\nabla^Hu$ is with respect to $g_H$) we deduce that $\langle(\nabla f)',\nabla^Hu\rangle=\mu^2\nu^2\langle\nabla^Hf,\nabla^Hu\rangle.$ So we have
$$(\nabla f)'=\nabla^Hf+\mu^4\nu^2\langle\nabla^Hf,\nabla^Hu\rangle\nabla^Hu$$
and deduce that 
\begin{equation}\label{pf lem grad nu eqn nabla f nabla H}
\nabla f=\nabla^Hf+\mu^2\nu^2\langle\nabla^Hf,\nabla^Hu\rangle\left(\mu^2\nabla^Hu+e_{n+1}\right)=\nabla^Hf+\mu^2\nu\langle\nabla^Hf,\nabla^Hu\rangle N.
\end{equation}
Using (\ref{pf lem grad nu eqn nabla f nabla H}) for $f=u$ and $\mu$, and computing using (\ref{notation est C1 expr N}) with $\nabla^Su=\mu^2\nabla^Hu$ we easily obtain
$$\langle\nabla u,\nabla\mu\rangle=\langle\nabla^H u,\nabla^H\mu\rangle\left(1+\mu^4\nu^2|\nabla^Hu|^2\right),$$
which implies (\ref{pf lem grad nu eqn interm}) and ends the proof of the lemma.
\end{proof}

\end{document}
